% GENERATED FILE. Edit canonical lessons, then run npm run generate:books.
% canonical-source-hash: fnv1a64:c1d8160718f7fb3d
% canonical-lessons: RU-C94-television, RU-C94-computer, RU-C94-fridge, RU-C94-cooker, RU-C94-shower

\chapter{Television, Computer, Fridge, Cooker, Shower}
\label{ch:ru-television-computer-fridge-cooker-shower}
\hlchaptermodality{94}

\begin{chapteropening}
\textbf{By the end of this chapter:} \emph{I can say a television, a computer, a fridge, a cooker and a shower.}

Complete the last lesson of chapter 94: I can say a television, a computer, a fridge, a cooker and a shower.
\end{chapteropening}

\section[\texorpdfstring{\ru{телевизор} (televízor)}{телевизор (televízor)}]{\ru{телевизор} (televízor) — a television}

\label{lesson:RU-C94-television}



\begin{quote}
Before the new one: say the Russian for holidays, leave, then the Russian for a rug, a carpet.
\end{quote}

\subsection*{You'll want to know: \ru{телевизор}}
\textbf{\ru{телевизор}} (\emph{televízor}) — \textquotedblleft{}a television\textquotedblright{}.

\begin{cousinweb}
\textbf{\ru{телевизор}} is an international coinage: Greek \emph{tēle}, \textquotedblleft{}far\textquotedblright{}, with \emph{-\ru{визор}}, from Latin \emph{vīsor}, \textquotedblleft{}one who sees\textquotedblright{}, from \emph{vidēre}, \textquotedblleft{}to see\textquotedblright{}: a far-seer.
\end{cousinweb}

\begin{grammarlens}[title={What you've built}]
One more word for this chapter.
\end{grammarlens}

\subsection*{Your turn}
\begin{itemize}
\raggedright
  \item \emph{Say it:} \emph{televízor}
  \item \emph{Say it:} \emph{televízor}, once more
  \item \emph{Say it:} say \emph{televízor}
  \item \emph{Recall:} say the Russian for holidays, leave, then the Russian for a rug, a carpet, then say \emph{televízor} again
\end{itemize}

\begin{tcolorbox}[breakable,colback=teal!4,colframe=teal!35!black,title={Before you move on}]
What does \ru{телевизор} mean? (\textquotedblleft{}A television\textquotedblright{}.) Say it once more.
\end{tcolorbox}
\section[\texorpdfstring{\ru{компьютер} (kompyúter)}{компьютер (kompyúter)}]{\ru{компьютер} (kompyúter) — a computer}

\label{lesson:RU-C94-computer}



\begin{quote}
Before the new one: say the Russian for a rug, a carpet, then the Russian for a television.
\end{quote}

\subsection*{You'll want to know: \ru{компьютер}}
\textbf{\ru{компьютер}} (\emph{kompyúter}) — \textquotedblleft{}a computer\textquotedblright{}.

\begin{cousinweb}
\textbf{\ru{компьютер}} comes from English \emph{computer}, from Latin \emph{computāre}, \textquotedblleft{}to reckon\textquotedblright{}.
\end{cousinweb}

\begin{grammarlens}[title={What you've built}]
One more word for this chapter.
\end{grammarlens}

\subsection*{Your turn}
\begin{itemize}
\raggedright
  \item \emph{Say it:} \emph{kompyúter}
  \item \emph{Say it:} \emph{kompyúter}, once more
  \item \emph{Say it:} say \emph{kompyúter}
  \item \emph{Recall:} say the Russian for a rug, a carpet, then the Russian for a television, then say \emph{kompyúter} again
\end{itemize}

\begin{tcolorbox}[breakable,colback=teal!4,colframe=teal!35!black,title={Before you move on}]
What does \ru{компьютер} mean? (\textquotedblleft{}A computer\textquotedblright{}.) Say it once more.
\end{tcolorbox}
\section[\texorpdfstring{\ru{холодильник} (kholodílnik)}{холодильник (kholodílnik)}]{\ru{холодильник} (kholodílnik) — a fridge}

\label{lesson:RU-C94-fridge}



\begin{quote}
Before the new one: say the Russian for a television, then the Russian for a computer.
\end{quote}

\subsection*{You'll want to know: \ru{холодильник}}
\textbf{\ru{холодильник}} (\emph{kholodílnik}) — \textquotedblleft{}a fridge\textquotedblright{}.

\begin{cousinweb}
\textbf{\ru{холодильник}} is \emph{\ru{холодить}}, \textquotedblleft{}to chill\textquotedblright{}, from \emph{\ru{холод}}, \textquotedblleft{}cold\textquotedblright{}, with \emph{-\ru{льник}}, a suffix for a tool: \textquotedblleft{}a chiller\textquotedblright{}.
\end{cousinweb}

\begin{grammarlens}[title={What you've built}]
One more word for this chapter.
\end{grammarlens}

\subsection*{Your turn}
\begin{itemize}
\raggedright
  \item \emph{Say it:} \emph{kholodílnik}
  \item \emph{Say it:} \emph{kholodílnik}, once more
  \item \emph{Say it:} say \emph{kholodílnik}
  \item \emph{Recall:} say the Russian for a television, then the Russian for a computer, then say \emph{kholodílnik} again
\end{itemize}

\begin{tcolorbox}[breakable,colback=teal!4,colframe=teal!35!black,title={Before you move on}]
What does \ru{холодильник} mean? (\textquotedblleft{}A fridge\textquotedblright{}.) Say it once more.
\end{tcolorbox}
\section[\texorpdfstring{\ru{плита} (plitá)}{плита (plitá)}]{\ru{плита} (plitá) — a cooker}

\label{lesson:RU-C94-cooker}



\begin{quote}
Before the new one: say the Russian for a computer, then the Russian for a fridge.
\end{quote}

\subsection*{You'll want to know: \ru{плита}}
\textbf{\ru{плита}} (\emph{plitá}) — \textquotedblleft{}a cooker\textquotedblright{}.

\begin{cousinweb}
\textbf{\ru{плита}} — Old Russian \emph{\ru{плита}}, \textquotedblleft{}a stone slab, a brick\textquotedblright{}. Vasmer thinks it most likely a cousin of Greek \emph{plínthos}, \textquotedblleft{}brick\textquotedblright{}; the common view takes it as a loan from that Greek word. A \emph{\ru{плита}} is the slab you cook on.
\end{cousinweb}

\begin{grammarlens}[title={What you've built}]
One more word for this chapter.
\end{grammarlens}

\subsection*{Your turn}
\begin{itemize}
\raggedright
  \item \emph{Say it:} \emph{plitá}
  \item \emph{Say it:} \emph{plitá}, once more
  \item \emph{Say it:} say \emph{plitá}
  \item \emph{Recall:} say the Russian for a computer, then the Russian for a fridge, then say \emph{plitá} again
\end{itemize}

\begin{tcolorbox}[breakable,colback=teal!4,colframe=teal!35!black,title={Before you move on}]
What does \ru{плита} mean? (\textquotedblleft{}A cooker\textquotedblright{}.) Say it once more.
\end{tcolorbox}
\section[\texorpdfstring{\ru{душ} (dush)}{душ (dush)}]{\ru{душ} (dush) — a shower}

\label{lesson:RU-C94-shower}



\begin{quote}
Before the new one: say the Russian for a fridge, then the Russian for a cooker.
\end{quote}

\subsection*{You'll want to know: \ru{душ}}
\textbf{\ru{душ}} (\emph{dush}) — \textquotedblleft{}a shower\textquotedblright{}.

\begin{cousinweb}
\textbf{\ru{душ}} comes from French \emph{douche}, from Italian \emph{doccia}, \textquotedblleft{}a pipe, a spout\textquotedblright{} (Vasmer).
\end{cousinweb}

\begin{grammarlens}[title={What you've built}]
One more word for this chapter.
\end{grammarlens}

\subsection*{Your turn}
\begin{itemize}
\raggedright
  \item \emph{Say it:} \emph{dush}
  \item \emph{Say it:} \emph{dush}, once more
  \item \emph{Say it:} say \emph{dush}
  \item \emph{Recall:} say the Russian for a fridge, then the Russian for a cooker, then say \emph{dush} again
\end{itemize}

\begin{tcolorbox}[breakable,colback=teal!4,colframe=teal!35!black,title={Before you move on}]
What does \ru{душ} mean? (\textquotedblleft{}A shower\textquotedblright{}.) Say it once more.
\end{tcolorbox}
